We approximated a solution to a modified Lotka-Volterra system of differential equations with several different numerical methods. We assessed the error via comparison to a highly converged RKF45 solution, and were able to achieve errors on the order of $10^{-2}$ with every method, except Euler's. Of the methods we implemented, we found RK4 to be the most efficient in terms of average step size. 

We were able to dictate which species survived by varying only $\beta_2$, and found a threshold value that produced the illusion of coexistence to binary precision. We derived an inequality system (Eq. ~\ref{eq:stability-conditions}) which guarantees coexistence, then satisfied it through variation in $\beta_1$. We also found that the introduction of a predator allowed coexistence under our initial parameterization due to the limiting of competition, whereas in the absence of a predator, one prey species became extinct. This result is consistent with the broader ecological principle that a shared predator can mediate competition, effectively acting as a stabilizing force on prey diversity.

Comparison to experimental and biological findings revealed several limitations of our model. Several parameters, particularly carrying capacities and competition coefficients, would realistically be interdependent in an ecological system, limiting the quantitative applicability of our results.
Future work could incorporate dynamic capacities or prey responses to predation to model the system more accurately. 
Regardless, the qualitative dynamics presented here align with published ecological findings and demonstrate the utility of numerical methods in exploring systems without analytical solutions. 